\documentclass[11pt,a4paper]{article}
\usepackage[utf8]{inputenc}
\usepackage[T1]{fontenc}
\usepackage[french]{babel}
\usepackage{geometry}
\geometry{margin=2.3cm}
\usepackage{booktabs}
\usepackage{longtable}
\usepackage{array}
\usepackage{amsmath}
\usepackage{xcolor}
\usepackage{hyperref}

\title{Parametres Unstructured KB: Strategie Courante, Variantes et Choix Final}
\author{Projet PFE STM32Cube}
\date{\today}

\begin{document}
\maketitle

\section{Objectif du document}

Ce document est \textbf{independant} et explique:
\begin{itemize}
  \item la configuration courante de la base KB Unstructured,
  \item les differentes strategies candidates,
  \item les differences de comportement entre strategies,
  \item la methode de comparaison pour retenir une strategie finale.
\end{itemize}

\section{Configuration courante (Current)}

\begin{longtable}{p{5cm}p{9cm}}
\toprule
\textbf{Parametre} & \textbf{Valeur courante} \\
\midrule
Search type & Hybrid \\
Semantic threshold & 0.55 \\
Semantic max doc. count & 25 \\
Full-text max doc. count & 12 \\
\bottomrule
\end{longtable}

\subsection{Signification des parametres}

\begin{itemize}
  \item \textbf{Search type = Hybrid}: combine la recherche semantique (proximite de sens) et la recherche full-text (correspondance exacte de termes techniques).
  \item \textbf{Semantic threshold = 0.55}: score minimum pour accepter un document dans la branche semantique.
  \item \textbf{Semantic max doc. count = 25}: limite superieure des documents semantiques injectes au contexte.
  \item \textbf{Full-text max doc. count = 12}: limite superieure des documents full-text injectes au contexte.
\end{itemize}

\section{Effet de chaque levier}

\subsection{Search type}
\begin{itemize}
  \item \textbf{Hybrid}: robuste sur questions mixtes (langage naturel + noms d'API/fichiers).
  \item \textbf{Semantic only}: meilleur sur paraphrases, plus fragile sur identifiants exacts.
  \item \textbf{Full-text only}: excellent sur correspondances exactes, faible sur reformulations.
\end{itemize}

\subsection{Semantic threshold}
\begin{itemize}
  \item \textbf{Seuil plus bas} (ex: 0.45): augmente le rappel, mais ajoute du bruit.
  \item \textbf{Seuil plus haut} (ex: 0.65): augmente la precision, mais peut rater des documents utiles.
\end{itemize}

\subsection{Nombre max de documents}
\begin{itemize}
  \item \textbf{Plus de docs semantiques}: meilleure couverture de sens, risque de dilution du contexte.
  \item \textbf{Plus de docs full-text}: meilleure couverture des tokens exacts (HAL\_, LL\_, fichiers), risque de redondance.
\end{itemize}

\section{Strategie courante vs strategies alternatives}

\begin{longtable}{>{\raggedright\arraybackslash}p{3.4cm}c c c c >{\raggedright\arraybackslash}p{5.2cm}}
\toprule
\textbf{Strategie} & \textbf{Search} & \textbf{Sem. thr.} & \textbf{Sem. docs} & \textbf{FT docs} & \textbf{Ce que le changement fait concretement} \\
\midrule
S0 Current (Balanced) & Hybrid & 0.55 & 25 & 12 & Point d'equilibre actuel entre precision et rappel. \\
S1 Recall+ & Hybrid & 0.45 & 35 & 15 & Elargit fortement la collecte; capte plus de cas marginaux mais augmente le bruit. \\
S2 Precision+ & Hybrid & 0.65 & 18 & 8 & Filtre agressif; moins de bruit mais risque de faux negatifs sur questions implicites. \\
S3 Semantic-heavy & Hybrid & 0.55 & 35 & 6 & Favorise le sens global; moins d'ancrage exact sur noms d'API/fichiers. \\
S4 Fulltext-heavy & Hybrid & 0.55 & 15 & 20 & Favorise les correspondances exactes; moins robuste aux reformulations. \\
S5 Semantic-only & Semantic & 0.55 & 25 & 0 & Supprime l'appui full-text; utile pour ablation, moins robuste en support STM32. \\
S6 Fulltext-only & Full-text & -- & 0 & 20 & Supprime la branche semantique; performant sur identifiants exacts, faible sur langage naturel. \\
\bottomrule
\end{longtable}

\section{Methode de choix final}

La strategie finale est selectionnee sur un protocole unique, applique \textbf{a toutes les strategies} sur le meme jeu de questions:
\begin{itemize}
  \item Quality Index global,
  \item pertinence Top-1/Top-3 des documents,
  \item taux d'hallucination,
  \item qualite des citations (verifiabilite),
  \item latence moyenne de reponse.
\end{itemize}

\subsection{Regle de decision recommandee}

On retient la strategie qui maximise le Quality Index \textbf{sous contraintes}:
\begin{equation}
\text{Choix final} = \arg\max_s\; QI(s)
\quad \text{s.c.} \quad
\begin{cases}
Hallucination(s) \leq H_{max}\\
Latence(s) \leq L_{max}\\
Citation\_Quality(s) \geq C_{min}
\end{cases}
\end{equation}

\section{Pourquoi la configuration courante est defendable}

La configuration \textbf{Hybrid, 0.55, 25, 12} est defendable car:
\begin{itemize}
  \item elle limite le bruit semantique tout en gardant les matches exacts,
  \item elle couvre le double besoin STM32: concepts + identifiants techniques,
  \item elle reste stable operationnellement sans saturation excessive du contexte.
\end{itemize}

\section{Message synthetique pour encadrantes}

\textbf{En une phrase}: la strategie courante est un compromis pragmatique entre rappel, precision et robustesse technique; les variantes servent a quantifier les risques de sur-filtrage (precision+) ou de bruit (recall+), puis a justifier objectivement le choix final par tableau comparatif et metriques communes.

\end{document}
